% ===== FILE 5/9: Split nguyên bản, Sigmoid nghịch đảo, Quỹ đạo opacity reset, Histogram reset, Tổng kết pipeline gốc =====

% --- Slide: Split (3DGS nguyên bản) ---
\begin{frame}{Split nguyên bản: lấy mẫu ngẫu nhiên quanh Gaussian cha}
\scriptsize
Khi một Gaussian \textbf{to} (trục dài nhất $>$ ngưỡng \texttt{percent\_dense}$\times$extent) vẫn còn gradient trị
tuyệt đối cao, cơ chế \textbf{split nguyên bản} (\texttt{densify\_and\_split}, \texttt{gaussian\_model.py:866--889})
sinh ra $N=2$ Gaussian con để "chia việc" khớp chi tiết:
\[
x_{\text{con}} = \mu_{\text{cha}} + R(q)\,\varepsilon,\qquad \varepsilon\sim\mathcal{N}(0,\mathrm{diag}(s_{\text{cha}})^2),
\qquad s_{\text{con}} = \frac{s_{\text{cha}}}{0.8\,N} = \frac{s_{\text{cha}}}{1.6}
\]
\begin{itemize}\itemsep1pt
    \item Vị trí mỗi con được \textbf{lấy mẫu ngẫu nhiên} từ một phân phối Gaussian có độ lệch chuẩn chính bằng scale
    của cha, sau đó xoay theo hướng $R(q)$ của cha --- tức là các con rải rác quanh cha, theo đúng hình dạng ellipsoid
    của cha (không rải đều hình cầu).
    \item Cả $N=2$ con đều bị chia scale cho \textbf{cùng một hệ số $1.6$ trên cả 3 trục}, bất kể trục nào mới thực sự
    là trục "quá to" gây ra vấn đề --- đây là điểm mà cơ chế split dị hướng của SADGS (phần sau bài giảng) cải tiến.
    \item Gaussian cha bị xoá ngay sau khi sinh $2$ con --- tổng số Gaussian tăng thêm đúng $1$ đơn vị mỗi lần một điểm
    bị split.
\end{itemize}
\begin{center}
\includegraphics[width=0.4\linewidth]{fig_09_split_geometry.png}
\end{center}
\end{frame}

% --- Slide: Sigmoid nghịch đảo & opacity reset ---
\begin{frame}{Vì sao lưu opacity dưới dạng logit? Và cơ chế reset định kỳ}
\scriptsize
Opacity $\alpha\in(0,1)$ được lưu nội bộ dưới dạng \textbf{logit} $o=\sigma^{-1}(\alpha)=\ln\!\frac{\alpha}{1-\alpha}$
(giá trị thực bất kỳ), rồi đi qua sigmoid $\alpha=\sigma(o)$ khi cần dùng --- giống hệt lý do dùng $\exp$ cho scale: để
mọi số thực đều tương ứng với một giá trị $\alpha$ \textbf{hợp lệ}, không cần ràng buộc \texttt{clamp} thủ công trong
quá trình tối ưu bằng gradient descent.
\begin{itemize}\itemsep1pt
    \item \textbf{Vì sao cần reset opacity định kỳ?} Nếu để tự nhiên, nhiều Gaussian "chết dần" (opacity giảm về rất
    thấp) vẫn tồn tại lâu, chiếm bộ nhớ mà không đóng góp gì đáng kể. Mỗi \texttt{opacity\_reset\_interval}$=3000$
    iteration (\emph{với điều kiện} $\texttt{iteration}<\texttt{densify\_until\_iter}=15\,000$, tức \textbf{chỉ đúng
    $4$ lần}: tại $3\text{k},6\text{k},9\text{k},12\text{k}$), toàn bộ opacity bị \textbf{giảm mạnh}:
    \[
    \alpha \leftarrow \alpha_{\text{cũ}} \times \texttt{opacity\_reset\_decay} = \alpha_{\text{cũ}}\times 0.1
    \]
    \item \textbf{Cơ chế "sàng lọc":} Gaussian nào thực sự hữu ích (đang đóng góp giảm loss) sẽ nhanh chóng \textbf{hồi
    phục} opacity qua vài trăm bước gradient tiếp theo; Gaussian nào không còn đóng góp thì opacity vẫn ở mức thấp và
    sẽ bị prune ở đợt dọn dẹp kế tiếp (slide sau).
    \item Vì reset \textbf{dừng hẳn} sau iteration $15\,000$, nửa sau của quá trình train không còn cơ chế chủ động nào
    "sàng lọc lại" opacity định kỳ nữa.
\end{itemize}
\begin{center}
\includegraphics[width=0.30\linewidth]{fig_16_sigmoid_inverse.png}
\end{center}
\end{frame}

% --- Slide: Quỹ đạo opacity qua các lần reset ---
\begin{frame}{Quan sát quỹ đạo opacity của một Gaussian qua nhiều lần reset}
\scriptsize
Theo dõi $\alpha(t)$ của một Gaussian cụ thể xuyên suốt quá trình train cho thấy rõ "nhịp điệu" của cơ chế reset: mỗi
lần chạm mốc reset ($3\text{k},6\text{k},9\text{k},12\text{k}$), $\alpha$ bị "đập" xuống còn $10\%$ giá trị trước đó,
sau đó \textbf{leo dần trở lại} nếu Adam vẫn còn nhận được gradient khuyến khích tăng opacity tại vị trí đó.
\begin{itemize}\itemsep1pt
    \item Một Gaussian \textbf{quan trọng} (đang giúp giảm loss thật sự): sau mỗi lần reset, $\alpha$ hồi phục nhanh,
    thường vượt lại mức trước khi reset trong vòng vài trăm iteration.
    \item Một Gaussian \textbf{dư thừa} (không còn đóng góp rõ rệt, ví dụ đã bị Gaussian khác lân cận "lấn sân"): sau
    reset, $\alpha$ hồi phục rất chậm hoặc gần như không hồi phục --- nó sẽ nằm dưới ngưỡng prune ở lần dọn tiếp theo.
    \item Cơ chế này về bản chất là một dạng \textbf{"kiểm tra sức khỏe" định kỳ} cho toàn bộ quần thể Gaussian, tách
    biệt hẳn với cơ chế clone/split/prune dựa trên gradient/kích thước đã nói ở các slide trước.
\end{itemize}
\begin{center}
\includegraphics[width=0.36\linewidth]{fig_17_opacity_trajectory.png}
\end{center}
\end{frame}

% --- Slide: Histogram trước/sau reset ---
\begin{frame}{Nhìn trên toàn bộ quần thể: phân bố opacity trước và sau reset}
\scriptsize
Thay vì theo dõi một Gaussian, hãy nhìn \textbf{toàn bộ phân bố} opacity của mọi Gaussian ngay trước và ngay sau một
lần reset. Ngay sau reset, toàn bộ phân bố bị "nén" mạnh về phía các giá trị thấp (do nhân với $0.1$) --- gần như mọi
Gaussian đều tạm thời "mờ đi" đồng loạt.
\begin{itemize}\itemsep1pt
    \item Điều thú vị xảy ra ở các iteration \textbf{sau} reset: phân bố dần \textbf{tách thành hai nhóm} --- một nhóm
    hồi phục nhanh trở lại vùng opacity cao (Gaussian hữu ích), một nhóm tiếp tục nằm lại vùng thấp (ứng viên prune).
    \item Đây chính là hiệu ứng "sàng lọc" đã mô tả ở hai slide trước, nhưng nhìn dưới góc độ \textbf{thống kê quần
    thể} thay vì một cá thể --- một cách trực quan để kiểm chứng cơ chế reset có thực sự đang "phân loại" Gaussian hữu
    ích/dư thừa hay không trong một lần chạy huấn luyện cụ thể.
\end{itemize}
\begin{center}
\includegraphics[width=0.36\linewidth]{fig_18_hist_reset.png}
\end{center}
\end{frame}

% --- Slide: Tổng kết pipeline 3DGS gốc ---
\begin{frame}{Tổng kết: pipeline 3DGS gốc, đầu-cuối}
\scriptsize
Gộp lại toàn bộ phần vừa trình bày thành một vòng lặp duy nhất:
\[
\text{SfM/COLMAP} \to \{\mu,\Sigma,\alpha,c\}_{i=1}^N \xrightarrow{\text{EWA}} \text{ellipse 2D}
\xrightarrow{\text{compact box}} \text{tile} \xrightarrow{\text{blend}} C(x)
\xrightarrow{\mathcal{L}} \nabla\theta \xrightarrow{\text{Adam}} \theta_{\text{mới}}
\]
lặp lại hàng chục nghìn lần, với riêng mỗi $100$ iteration lại chen thêm bước \textbf{ADC} (clone/split/prune) dựa trên
gradient tích luỹ, và mỗi $3000$ iteration (tới $15\,000$) lại reset opacity một lần.
\begin{itemize}\itemsep1pt
    \item \textbf{Điểm mạnh của thiết kế này:} hoàn toàn tường minh và khả vi từ đầu đến cuối --- không hộp đen, mọi
    bước đều có thể kiểm tra, gỡ lỗi; render thời gian thực ngay sau khi train xong.
    \item \textbf{Điểm còn hạn chế, là động lực cho phần tiếp theo của bài giảng:} tín hiệu densify \textbf{chỉ dựa vào
    gradient} --- một đại lượng phản ánh \emph{trạng thái hội tụ hiện tại} của loss, chứ không trực tiếp đo "ảnh này còn
    thiếu bao nhiêu \% chi tiết tần số cao chưa được biểu diễn, theo tỉ lệ nào". Hai Gaussian có cùng độ lớn gradient có
    thể đang thiếu hụt hình học ở hai \emph{quy mô} rất khác nhau (một cái thiếu chi tiết rất mịn, một cái thiếu cấu
    trúc thô hơn nhiều) --- gradient một mình không phân biệt được điều này.
\end{itemize}
\begin{center}
\includegraphics[width=0.34\linewidth]{samples_train.png}\\
\captionof{figure}{\footnotesize Ảnh render minh hoạ (định tính, không kèm số đo benchmark).}
\end{center}
\end{frame}
